% ============================================================
% UNIDAD 8: CONECTORES MILITARES Y D-SUB
% ============================================================
\clearpage
\section*{UNIDAD 8: CONECTORES}
\subsection*{Circulares militares, D-sub, contactos y accesorios}

\begin{recuadro}
\textbf{Objetivo de la unidad:} reconocer los conectores usados en aviónica, interpretar su
designación, armar conectores con contactos crimpados e insertados, y usar correctamente sus
herramientas y accesorios.
\end{recuadro}

% ------------------------------------------------------------
\section*{1. CONCEPTOS GENERALES}

\begin{itemize}
  \item \textbf{Pin y casquillo:} el contacto macho (\textit{pin}) y el hembra (\textit{socket}).
        Regla de seguridad: el lado que puede quedar \textbf{energizado} al desconectar debe tener
        \textbf{casquillos} (hembras), para que no queden pines con tensión expuestos.
  \item \textbf{Ficha y receptáculo:} la parte que va en el extremo del cable (\textit{plug}) y la
        que va fija en el equipo o en un panel (\textit{receptacle}).
  \item \textbf{Inserto y arreglo de contactos:} la pieza aislante que sostiene los contactos y su
        distribución (cantidad, tamaño y posición).
  \item \textbf{Polarización o clave:} ranuras y chavetas que impiden conectar dos conectores
        iguales cruzados. Los conectores militares tienen varias posiciones de clave (N normal, A,
        B, C\ldots).
  \item \textbf{Contactos de crimpar y de soldar;} contactos \textbf{removibles} (se insertan y
        extraen con herramientas) o fijos.
\end{itemize}

% ------------------------------------------------------------
\section*{2. CONECTORES D-SUB (SUBMINIATURA D)}

Tienen carcasa en forma de D (que asegura la polarización), una o varias filas de contactos y
tornillos de fijación. La norma militar es la \textbf{MIL-DTL-24308}. Son los conectores de la
mayoría de los equipos de aviónica de aviación general (Garmin, Avidyne, Dynon, entre otros),
generalmente en \textbf{alta densidad} y con contactos de crimpar.

\begin{center}
\small
\begin{tabular}{|c|c|c|}
\hline
\textbf{Tamaño de carcasa} & \textbf{Densidad estándar} & \textbf{Alta densidad} \\ \hline
E & DE-9 (9 contactos) & DE-15 (15) \\ \hline
A & DA-15 (15) & DA-26 (26) \\ \hline
B & DB-25 (25) & DB-44 (44) \\ \hline
C & DC-37 (37) & DC-62 (62) \\ \hline
D & DD-50 (50) & DD-78 (78) \\ \hline
\end{tabular}
\end{center}

\begin{itemize}
  \item Los contactos de crimpar se insertan desde atrás hasta que traban y se extraen con una
        herramienta de extracción que libera la traba.
  \item Los contactos de alta densidad son más pequeños y requieren su propia crimpadora o
        posicionador (el que indica el fabricante del equipo).
  \item \textbf{Numeración:} cada contacto tiene su número grabado en el inserto. Al mirar de frente
        el conector macho y el hembra, la numeración queda espejada. En los planos se indica si la
        vista es ``de frente'' (\textit{mating face}) o ``desde atrás'' (\textit{wire side}):
        confundirla es un error clásico.
  \item \textbf{Carcasas traseras} (\textit{backshells}) metálicas para fijar el mazo, dar alivio de
        tensión y conectar el blindaje a masa.
\end{itemize}

% ------------------------------------------------------------
\section*{3. CONECTORES CIRCULARES MILITARES}

\begin{center}
\small
\begin{tabular}{|p{0.22\linewidth}|p{0.30\linewidth}|p{0.36\linewidth}|}
\hline
\textbf{Norma} & \textbf{Acoplamiento} & \textbf{Características y uso} \\ \hline
MIL-DTL-38999 & Serie I y II: bayoneta; Serie III: roscado autotrabante; Serie IV: de traba
rápida & Alta densidad, sellado ambiental, muy usado en aeronaves modernas, sensores y equipos. \\
\hline
MIL-DTL-26482 & Bayoneta & Tamaño reducido, uso general en aviónica. \\ \hline
MIL-DTL-5015 & Roscado & Robusto, para alimentación y cargas grandes; aeronaves antiguas. \\
\hline
MIL-DTL-83723 & Bayoneta o roscado & Aplicaciones aeronáuticas ambientales. \\ \hline
\end{tabular}
\end{center}

\paragraph*{Cómo se lee una designación MIL-DTL-38999.} Ejemplo: \texttt{D38999/26WD35SN}
\begin{center}
\small
\begin{tabular}{|c|p{0.66\linewidth}|}
\hline
\texttt{D38999} & Norma MIL-DTL-38999 \\ \hline
\texttt{/26} & Tipo de pieza: ficha recta de la serie III (otros números identifican receptáculos de
pared, de brida, pasamuros, etc.) \\ \hline
\texttt{W} & Terminación superficial de la carcasa (por ejemplo, cadmio verde oliva) \\ \hline
\texttt{D} & Tamaño de carcasa (letra A a J) \\ \hline
\texttt{35} & Arreglo de inserto (cantidad y tamaño de contactos) \\ \hline
\texttt{S} & Contactos hembra (\textit{socket}); \texttt{P} para macho (\textit{pin}) \\ \hline
\texttt{N} & Posición de clave normal \\ \hline
\end{tabular}
\end{center}

\paragraph*{Tamaños de contactos.} Se identifican por número: 22D, 20, 16, 12 (cuanto menor el
número, mayor el contacto y la corriente que admite). Cada tamaño acepta un rango de secciones de
cable. También hay contactos coaxiales, de fibra óptica y de alta potencia.

% ------------------------------------------------------------
\section*{4. ARMADO DE UN CONECTOR CON CONTACTOS REMOVIBLES}
\begin{enumerate}
  \item Verificar en el plano el número de parte del conector, de los contactos y la asignación de
        cada cable a cada número de contacto (\textit{pinout}).
  \item Pasar por el cable, antes de crimpar, los accesorios traseros: backshell, alivio de
        tensión, termocontraíbles.
  \item Crimpar cada contacto (Unidad 7) e inspeccionarlo.
  \item Insertar cada contacto por la parte trasera del inserto con la \textbf{herramienta de
        inserción} (de plástico, familia M81969/14) hasta escuchar o sentir la traba. Tirar
        suavemente del cable para verificar que quedó trabado.
  \item Colocar \textbf{tapones de sellado} en todas las cavidades sin uso: si quedan abiertas, se
        pierde el sellado ambiental.
  \item Armar el backshell, conectar el blindaje y ajustar al torque indicado.
  \item Verificar \textbf{pin a pin} la continuidad y el aislamiento contra el plano.
\end{enumerate}
Para extraer un contacto se usa la \textbf{herramienta de extracción} del tamaño correcto,
introducida desde atrás (o desde el frente, según el conector) para liberar la traba. Usar la
herramienta equivocada daña el inserto, y entonces hay que cambiar el conector.

\begin{recuadro}
\textbf{Conectores en la aeronave.} Los conectores roscados en zonas de vibración se aseguran con
\textbf{alambre de frenar} o son autotrabantes; los de bayoneta, con su traba. Antes de conectar:
inspeccionar que no haya pines doblados, retraídos o contactos ``empujados'' hacia atrás, suciedad o
humedad. Nunca forzar un conector: si no entra, algo está mal (clave, pin doblado).
\end{recuadro}

% ------------------------------------------------------------
\section*{5. OTROS ELEMENTOS DE INTERCONEXIÓN}
\begin{itemize}
  \item \textbf{Conectores de bandeja} (ARINC 404 y ARINC 600): en los racks de aviónica de los aviones
        de transporte, con contactos de señal, de potencia y coaxiales.
  \item \textbf{Bloques de unión y de masa} (sistemas de módulos con contactos crimpados, familia
        M81714): permiten unir varios cables a un mismo punto (por ejemplo, todas las masas de un
        equipo) sin empalmes.
  \item \textbf{Regletas de terminales} con tornillos: para cables con terminales de ojal.
  \item \textbf{Disyuntores (\textit{circuit breakers})}, interruptores, relés y fusibles (Unidad 9).
\end{itemize}

% ------------------------------------------------------------
\section*{6. ACTIVIDADES DE LA UNIDAD}
\begin{enumerate}
  \item Interpretar las designaciones \texttt{D38999/20WB98PN} y \texttt{D38999/26WJ61SA}
        (usar una tabla de arreglos de insertos).
  \item Armar un conector D-sub de alta densidad de 26 contactos siguiendo un pinout dado.
  \item Insertar y extraer contactos de un conector MIL-DTL-38999 con las herramientas correctas.
  \item Identificar los conectores de un equipo de aviónica real y buscar su número de parte en el
        manual de instalación.
\end{enumerate}
